\documentclass[10pt,a4paper]{article}
\usepackage[T1]{fontenc}
\usepackage[utf8]{inputenc}
\usepackage{mathptmx}
\usepackage{courier}
\usepackage[english]{babel}
\usepackage{amsmath,amssymb}
\usepackage[margin=2.3cm]{geometry}
\usepackage{microtype}
\usepackage{booktabs,tabularx,array}
\usepackage{titlesec}
\usepackage{caption}
\usepackage[hidelinks]{hyperref}

\titleformat{\section}{\normalsize\bfseries\scshape}{\thesection.}{0.6em}{}
\titleformat{\subsection}{\normalsize\bfseries}{\thesubsection}{0.6em}{}
\titleformat{\subsubsection}{\normalsize\itshape}{\thesubsubsection}{0.6em}{}
\titlespacing*{\section}{0pt}{1.4ex plus .4ex}{0.8ex}
\titlespacing*{\subsection}{0pt}{1.2ex plus .3ex}{0.6ex}
\titlespacing*{\subsubsection}{0pt}{1ex plus .3ex}{0.4ex}
\setlength{\parskip}{0.35ex}
\renewcommand{\arraystretch}{1.15}
\newcolumntype{Y}{>{\raggedright\arraybackslash}X}
\captionsetup{font=small,labelfont=bf}

\newcommand{\M}{\bar{M}}

\title{\large\bfseries Does Underlying Structure Predict\\ Quantum Trainability and Dequantization?}
\author{\normalsize Charbel A. Dargham}
\date{}

\begin{document}
\maketitle

\begin{abstract}
\noindent Quantum trainability and dequantization has, till this day, been shown
post-mortem of the construction and execution of a quantum machine learning
algorithm, after alot of effort and resources were put into making it. we
wish to find a way to know apriori whether or not a machine learning problem
is, first trainable with quantum circuits, and second know if a classical
algorithm would beat it. We narrow our vision in this study to a subcase of
learning problems.
\end{abstract}

%=====================================================================
\section{The Question}

\begin{quote}
\textbf{The Main Question of This Endeavor.}
\emph{Can the underlying structure of a learning problem predict, before training,
whether structure-preserving quantum models of it will be trainable and
resist dequantization?}
\end{quote}

\subsection{Why Underlying Structure}

McClean et al.\ 2018 (Nat.\ Commun.\ 9, 4812) found \textbf{barren plateaus}:
in random, structure-free parameterised circuits, gradients shrink
exponentially with the number of qubits. The training landscape is flat, so
the optimiser has no direction to follow.

Building the problem's symmetry into the circuit restricts what the circuit
can explore, and the gradients stay large, we site the following.
Larocca et al.\ 2022, \emph{Group-invariant QML}, PRX Quantum 3,
030341: how symmetry enters quantum models.
Ragone et al.\ 2024, Nat.\ Commun.\ 15, 7172, with Fontana et al.\ 2024 as
an independent derivation: an exact formula for gradient size in terms of
the \textbf{dynamical Lie algebra (DLA)}. The DLA is the space of all
operations the circuit's generators can produce. A small DLA gives large
gradients; an exponentially large DLA gives vanishing gradients.
Cerezo et al.\ 2021, Nat.\ Commun.\ 12, 1791: observables that act on
only a few qubits keep gradients alive in shallow circuits. This is where
locality comes in.

Most notably Cerezo et al.\ (Nat.\ Commun.\ 16, 7907, 2025) argue that the same
small DLA that makes a model trainable also lets a classical computer imitate
it. That imitation is \textbf{dequantization}, and it means the model
gives no quantum advantage. Goh et al.,
\emph{$\mathfrak{g}$-sim}, arXiv:2308.01432, is the concrete
classical simulation method.

We decided to expand on this work and perform an experiment a priori of any
proving theorems to understand how the interplay between four decided
structural features of a learning problem and our desired prediction behaves.
We define the four kinds of structure:
\textbf{symmetry}, transformations that don't change the answer;
\textbf{locality}, which parts of the input act on each other directly;
\textbf{spectrum}, which patterns or modes of the input carry the answer;
\textbf{data complexity}, how hard the inputs are to describe classically.

\textbf{Only symmetry has a theorem that runs from the problem to the
trainability formula.} Ragone's formula takes three inputs: the DLA
$\mathfrak{g}$, the input state $\rho$, and the measured observable
$O$. For locality, spectrum and data complexity, that mapping is the project's
own working hypothesis, not a published result.

\subsection{The Plan}

\subsubsection*{Phase A: Build the dataset of problems, circuits and measurements.}

(i)~Collect problems: these are real learning problems, we'll discuss this
step in Section~2, we have harvested, and after cleaning the data 1368
problems were harvested categorized into 4 types (more about that in
Section~2).
(ii)~Build circuits from each problem. There are three kinds:
a \textbf{natural model}, built from the problem's structure,
for example by \textbf{twirling}: averaging each gate over the
symmetry group so the gate becomes symmetric;
a \textbf{generic control} that ignores the structure;
a \textbf{path model} with a dial $\lambda$ running from generic
(0) to fully structured (1).
The controls are there to show whether the structure \emph{causes} good
behaviour or only comes with it.
(iii)~\textbf{Fingerprint each circuit} from its description alone, with
no training. The main quantities are the size of the space the circuit can
reach (the DLA) and how that space splits into blocks.
(iv)~\textbf{Commit predictions before measuring.} Apply Ragone et al.\
2024 for trainability and Goh et al.\ 2023 for the cost of classical
simulation, and commit the predictions to git. The commit timestamp proves
they came first.
(v)~\textbf{Measure.} Record how much the output varies at random
parameters, what happens in training, and what classical simulation cost.

\subsubsection*{Phase B: build the predictor.}

This is a statistical model that goes from fingerprint to (trainable?, hard to
simulate?).
\textbf{Testing:} it is always tested on \textbf{problem
families it has never seen}, so that it cannot succeed by memorising a
family.
\textbf{Success:} the predictor counts as a success if it beats
the theoretical formulas where they don't apply, such as shallow circuits,
approximate symmetry, or cases outside the theorems' conditions.

%=====================================================================
\section{Phase A}

\subsection{Data Harvest}

\subsubsection{What does the dataset contain?}

At each row of our dataset, we have a problem, one binary learning task at a
fixed $n$: its samples, its TRAIN/TEST split, its label rule, its equation and
$\M$. Effectively, one problem is
\[
\mathcal{P}_n = (\mathcal{X}_n,\ \mathcal{Y}_n,\ \mu_n,\ \mathcal{E}_n,\ f^{*}),
\]
whose parts are given in Table~\ref{tab:record}.

\begin{table}[h]
\centering\small
\caption{The parts of a problem record.}\label{tab:record}
\begin{tabularx}{0.85\linewidth}{@{}lY@{}}
\toprule
Part & Meaning\\
\midrule
$\mathcal{X}_n$ & the inputs, on $n$ sites, with $n$ = number of qubits, from 4 to 12\\
$\mathcal{Y}_n$ & the labels, always $\pm1$\\
$\mu_n$ & the data distribution: the actual samples and their fixed 70/30 TRAIN/TEST split\\
$\mathcal{E}_n$ & the differential equation each input determines, written as one operator $M(x)$ per input\\
$f^{*}$ & the target, plus the sentence saying how the labels were made\\
\bottomrule
\end{tabularx}
\end{table}

\textbf{$\M$, the task operator.} It is built from TRAIN only. $M(x)$
feeds the circuit's encoding; $\M$ feeds its trainable generators.
\textbf{Provenance:} source, licence, citation, file hash, group,
and the ``modelling map'' (how a domain object was turned into $M(x)$).
\textbf{An origin tier:}
tier 1, REAL: real inputs and real labels;
tier 2, REAL\_INPUT\_COMPUTED\_LABEL: real inputs, labels computed
exactly from them;
tier 3, SYNTHETIC: both generated.

\textbf{The type of a problem is fixed by its operator, not by its domain}
(HARVEST.md \S3). The same domain can hide different mathematics:
MaxCut encodes like a graph, while 3-SAT has three-qubit terms that a graph
does not have (Table~\ref{tab:types}).

\begin{table}[h]
\centering\small
\caption{The four problem types.}\label{tab:types}
\begin{tabularx}{\linewidth}{@{}l Y Y Y@{}}
\toprule
Type & $M(x)$ & Equation & Do sites have an identity?\\
\midrule
GRAPH & Laplacian $L = B^{\top} W B$ & heat, $\dot x = -Mx$ & no\\
COVARIANCE & precision $Q = R^{-1}$, the inverse of the channel correlation matrix & Ornstein--Uhlenbeck, $dx = -Mx\,dt + \text{noise}$ & yes (channel C3 is C3 in every trial)\\
DIAGONAL & clause-penalty Hamiltonian, $Z$ terms only & Schr\"odinger, with $e^{-i\gamma H}$ as the encoding & no\\
QUANTUM & Hamiltonian with non-commuting terms; the input is its ground state & Schr\"odinger, $\dot\psi = -iM\psi$ & yes (chain sites)\\
\bottomrule
\end{tabularx}
\end{table}

\textbf{Equation checks} (HARVEST.md \S4). These are run on
every input. Failures are counted, never fixed, and the final count is 0
failures.
Graphs: $L$ is symmetric with zero row sums, and the number of zero
eigenvalues equals the number of connected components.
Covariance: $QR = I$.
QUANTUM: $H\psi = E\psi$, and $E$ is the lowest eigenvalue.
SAT: the Pauli expansion equals the violated-clause count on all $2^n$
assignments.

The audit script ran on the raw data and found \textbf{1\,379 problems and 0
structural failures}. Files, counts, splits, labels, equation checks, finite
inputs, provenance and unique ids were all fine. The problems were in the
content (Table~\ref{tab:findings}).

\begin{table}[h]
\centering\small
\caption{Findings of the audit of the raw harvest.}\label{tab:findings}
\begin{tabularx}{\linewidth}{@{}c Y Y Y@{}}
\toprule
\# & Finding & Size & Cause\\
\midrule
1 & \textbf{Leak:} identical inputs on both sides of TRAIN/TEST & 22 problems (42 contain duplicates at all). Worst: NAE-3 $n=4$, where 71 of 90 TEST inputs are also in TRAIN & random formulas repeat at small $n$; mirror-image molecules give the same graph; EigenWorms stores repeated recordings\\
2 & \textbf{Class balance below the floor} & 1 problem: XOR-3 $n=10$ count-high, overall 0.19, TEST 0.14 & the median-tie rule only checked TRAIN\\
3 & \textbf{Degenerate ground states:} the stored vector is an arbitrary point of the ground space & 10\,005 of 41\,920 spin samples (24\%), in 127 of 262 problems & physical: odd-$n$ chains that conserve total spin, ordered phases, cluster-Ising edge states\\
4 & \textbf{Hartree--Fock did not converge}, so the labels are unreliable & 2 H$_2$O geometries, each present in 2 problems & numerical\\
5 & \textbf{Uninformative operator:} Ledoit--Wolf shrinkage $\geq 0.99$, so the operator is the identity matrix & 475 UEA samples (LSST 183, RacketSports 292) & 30--36 time points with bursty signals\\
\bottomrule
\end{tabularx}
\end{table}

\textbf{Raw data were never modified.} The final dataset was built from
the raw harvest by a script with six stated rules (R1--R6), and the raw data's
sha256 is identical before and after the build.

\textbf{Problems found in the raw harvest} (1\,379 problems, 0 structural
failures):
(1)~\textbf{Leak:} in 22 problems, identical inputs sat on both sides
of TRAIN/TEST.
(2)~\textbf{Class balance:} 1 problem was below the 0.2 floor.
(3)~\textbf{Degenerate ground states:} 24\% of spin-chain samples, so
the stored state was arbitrary.
(4)~\textbf{Unreliable labels:} 2 molecule geometries where
Hartree--Fock did not converge.
(5)~\textbf{Uninformative operators:} 475 sensor samples whose
operator is the identity matrix, which acts as a label shortcut.

\textbf{What was done:}
\textbf{Spin chains re-harvested} with ground states computed
inside a fixed symmetry sector (fixes 3). Samples with gap $<10^{-6}$ fell from
642 to 10 on a dry run.
\textbf{R1:} excluded the unbalanced problem (fixes 2).
\textbf{R2:} group-aware re-split, keeping identical inputs on one
side, then recomputing labels and $\M$ on the new TRAIN (fixes 1).
\textbf{R3:} weights so each distinct input counts once in TEST.
\textbf{R4:} at least 30 distinct TEST inputs.
\textbf{R5:} removed the non-converged samples (fixes 4).
\textbf{R6:} flagged rather than removed the 475 uninformative and
26 degenerate samples (handles 5). A sensitivity check runs with and without
them.

\textbf{Final state:}
\textbf{1\,368 problems, 58 groups, 337\,930 samples.}
\textbf{Audit: 0 FAIL.} TRAIN/TEST leaks went from 22 problems to
0.
\textbf{2 WARN:} two TEST balances at 0.17, handled by reporting
balanced accuracy.
\textbf{One known gap:} XY-periodic (10 problems) is missing because a
solver failure made the harvest reject the whole dataset. The bug is fixed,
but the re-harvest is still pending.

%---------------------------------------------------------------------
\subsubsection{Do the four types differ in structure?}

Take one structure feature at a time. Once size $n$ and source are accounted
for, does its value depend on the type?

It tests one way the ``extract structure'' step could fail. failure as
features that are ``identical everywhere, or different only by construction''.
There are two possibilities.
\textbf{Features identical across types} carry no information, and
could never predict anything about circuits.
\textbf{Features that differ only because of how a type is
defined} are not evidence of anything. For example, every COVARIANCE
precision matrix is dense, so its support density is $\approx 1$ by
definition. A test that ``finds'' this has found nothing.

So Q2 asks whether the extracted structure varies \emph{beyond what we built
into it}. we decided to fit each feature into a hierarchicha generalized
addtitive mixed model (Wood 2017; Pedersen et al.\ 2019)
\[
\text{feature} = \text{type baseline} + \text{smooth curve in } n \text{ for each type} + \text{offset for each group} + \text{noise}.
\]
a B-spline, i.e.\ short polynomial pieces joined smoothly, limited to 3
degrees of freedom. That is enough to follow a curve like $2/n$ but not enough
to chase noise. Each source gets its own baseline, drawn from a common
distribution. This handles the fact that problems from one source are not
independent. Treating them as independent can push the false-positive rate
from 0.05 up to 0.80 (Aarts et al.\ 2014).

We effectively can read the following off the model.
\textbf{ICC (intraclass correlation)} = between-group variance /
total variance. Converted into an \textbf{effective sample size} (Kish):
how many independent observations the feature really has.
\textbf{Level test:} do the types differ at equal $n$? This is a
likelihood-ratio test of the model with type terms against the model without
them.
\textbf{Shape test:} do the types scale differently with $n$? One
curve per type against one shared curve.
\textbf{Holm correction} over the 12 features. It keeps the chance
of even one false ``difference'' at 5\%.
\textbf{Marginal $R^2$} (Nakagawa \& Schielzeth
2013): the share of variance explained by type and $n$ alone, without the
group offsets.
\textbf{Pairwise type differences} at $n = 8$, 10 and 12.

Every forced value matches its derivation exactly
(\texttt{q2\_pipeline\_check} is empty). One bug was found this way and
fixed before any result was read: factor labels depended on site order.

\textbf{The spectral features are the clear finding.} Spectral
entropy, participation ratio and algebraic connectivity:
\textbf{level:} they differ by type, with $p$ from $10^{-12}$ to $5\cdot10^{-7}$;
\textbf{shape:} for entropy and participation ratio, the types
also scale differently with $n$ (shape $p \approx 10^{-23}$, $2\cdot10^{-28}$);
\textbf{effect size:} marginal $R^2$ of 0.30--0.55.
\textbf{GRAPH stands apart.} It differs from every other type on all
three spectral features (pairwise $p < 10^{-4}$). COVARIANCE, DIAGONAL and QUANTUM
are mostly indistinguishable from each other (spectral entropy at $n=10$: 0.973,
0.963, 0.979; pairwise $p$ 0.21--0.66).
\textbf{Locality features} (support density, site-degree spread,
terms per site, symmetry classes) differ in level but not in shape. Their
``QUANTUM'' finding rows are only the 6 molecules at $n = 8$ and 12, because the
spin-chain cells are forced.
\textbf{Term-size fractions and the symmetry fraction show no
difference} after Holm correction ($p$ 0.08--1).
\textbf{ICC is 0.5--0.9}, so the effective sample size is
\textbf{18--132, not 1\,107.} The group-level model was necessary.
\textbf{Sensitivity:} results are identical with and without the
flagged samples.

The implications below are from ANALYSIS\_RESULTS ``What the study supports''
\S2 and \S5.
(1)~\textbf{The extracted structure varies beyond construction}, most
robustly through the weighted spectral features. It is not a relabelling of
the definitions.
(2)~\textbf{The variation is coarse.} It separates graphs from
everything else, but not covariance, SAT and quantum Hamiltonians from each
other. A predictor built on these features will see roughly two kinds of
operator, not four.
(3)~\textbf{Sources dominate.} Because the effective sample size is
small, the circuit phase should favour diversity of sources over more
problems per source, and should validate by holding out whole sources.

%---------------------------------------------------------------------
\subsubsection{Can a problem's kind be recovered from its structure alone?}

if the features only re-label the type, they carry nothing beyond the
definitions. We performed a logistic regression from the features to the
class. to valitdate we trained on all sources but one, test on the held-out
one and output a score for balanced accuracy.
\textbf{Three size controls:}
an $n$-only classifier as a baseline;
the Dinga et al.\ 2020 test, which asks whether the predictions add
information beyond $n$;
a subset balanced for size.
\textbf{Two layers:}
\textbf{Q3a} recovers the 4 types. This is a pipeline check.
\textbf{Q3b} recovers the domain within a type. This is the
finding.

as for the results:
\textbf{Q3a: 0.91} (chance 0.25, $n$-only 0.39). The errors fall
where the operator definitions predict. For example, NAE-3 and 2-SAT have
only pairwise terms and are classed as GRAPH.
\textbf{Q3b is negative.} COVARIANCE scores 0.49 (chance 0.50).
GRAPH (0.08) and DIAGONAL (0.20) are below chance, a known effect of
leave-one-group-out with few sources (Jamalabadi et al.\ 2016). QUANTUM's
1.00 is trivial: forced spin-chain features separate it.

the features identify the kind of operator, not the domain within a type.
There are two open readings:
the features are too coarse, because the factor graph loses how terms
overlap and the Hodge complex would keep it;
or the domain is not the partition that matters for trainability.

%---------------------------------------------------------------------
\subsubsection{How does structure change with size $n$?}

A predictor trained at $n \leq 12$ must extrapolate, and a feature with no
regular dependence on $n$ cannot be extrapolated. Each sub-family is followed
across $n$ as one line (43 of 64 lines have enough sizes), and each feature is
read in three tiers: the slope against $\log n$ with its 95\% interval; a
test of laws known in advance, exact for forced laws and by likelihood ratio
for predictable ones; and, where no law is known, a choice among constant,
$a+b/n$, $a+b\log n$, power and exponential, named only if
$\Delta\mathrm{BIC} > 10$ (Kass \& Raftery 1995). On the spin chains, two
physics laws are tested on each side of every transition: ground-state
entanglement grows as $(c/6)\log n$ on open chains and $(c/3)\log n$ on
periodic ones at criticality (Calabrese \& Cardy 2004) and stays bounded where
the chain is gapped (Hastings 2007); the in-sector gap falls as $1/n$ at
criticality.

The structure changes monotonically almost everywhere: 379 of 428 trends are
monotone, spectral entropy rises on 41 of 43 lines and support density falls
on all 39. All 28 forced laws match exactly. Among the predictable SAT laws,
PLANTED-3 agrees with its prediction, while NAE-3 deviates in a way explained
by its pairwise terms cancelling. Only 19 of 177 laws can be named, mostly on
the EEG lines; the other 158 are undetermined on $n = 4$--12. On the open XXZ
chains in their critical phase, the entanglement slope is 0.18--0.21 against
the predicted $c/6 = 0.167$, inside the interval in all four cases. Four
readings become flat when the flagged degenerate ground states are removed,
and are reported as depending on them.

The extracted structure therefore changes lawfully with $n$ and reproduces
known physics, but most laws cannot be identified at $n \leq 12$: a predictor
trained on these data describes finite sizes, not asymptotics.

%=====================================================================
\section{Conclusions}

The study validates the second link of the chain, the extracted structure,
before the expensive steps that follow; it says nothing directly about
trainability or dequantization, which require simulated circuits. Within that
scope, the extraction is correct: every forced value and every forced law
matches exactly, the one bug found by these checks was fixed before any result
was read, and the type-classification errors fall where the operator
definitions predict. The structure is varied beyond construction, most
robustly through the weighted spectral features, though coarsely: GRAPH
separates from the other three types, while COVARIANCE, DIAGONAL and QUANTUM
remain close. It changes lawfully with $n$, with the critical XXZ chains
following the Calabrese--Cardy law. It identifies the kind of operator but
not the domain within a type. Finally, the data hold far fewer independent
observations than problems (ICC 0.5--0.9, effective sample size 18--132), so
the circuit phase should favour diversity of sources and validate the
predictor by holding out whole sources.

%=====================================================================
\section{Next Steps}

The permutation tests of Q2 (5\,000 group-level shuffles per feature) and Q3
(1\,000 per task) remain to be run; until then the $p$-values rest on the
mixed models' assumptions. The robustness study (Q6) was not carried out, so
robustness currently rests on the R6 sensitivity runs and the size controls of
Q3. The analysis choices left open (classifier regularisation, the
correlation-block threshold, the minimum number of sizes per line, log
transforms) must be confirmed. On the data side, the XY-periodic spin chains
are to be re-harvested with the corrected solver, and the origin tier of the
molecule problems aligned across files. On the structure side, Hodge features
built on the factor graph would keep how terms overlap, which the current
features lose, and are the first candidate to test the ``features too
coarse'' reading of Q3b. The project then moves to the rest of Phase~A:
building natural, generic and path circuits from each problem, fingerprinting
them, committing the theoretical predictions, and measuring trainability and
the cost of classical simulation.

\end{document}
